\section{问题背景和分析}

\subsection{问题背景}
无线电干扰源的搜索与清除同时包含主动感知和路径规划：机器人需要移动到合适位置获取方向信息，又要控制移动、检测、频道切换和清除产生的时间成本。由于目标数量、位置、有效接收半径和发射方向均不可直接观测，单次测量不足以确定目标；同时，错误退出会遗漏未发现频道，因此策略既要追求效率，也要给出发现、清除和退出的可靠依据。

本文将四问组织为一条递进技术路线：问题一构造有界误差下的定位与清除证书，问题二选择兼顾观测可靠性和定位质量的第二检测点，问题三建立全向多源的完备搜索骨架和频道状态机，问题四把覆盖保证推广到含未知发射方向的定向源场景。

\subsection{问题一分析}
每次测向形成一个朝前的角度扇形。多个扇形相交后得到干扰源可行域。问题一除要求计算区域直径外，还需判断由最远点对构造的直径圆能否覆盖该区域；为此必须区分“区域直径”和“覆盖区域所需的最小半径”。

\subsection{问题二分析}
第二检测点既要与已有测向方向形成良好交会角，又要避免因未知接收半径而失去信号。本文先限制候选点位于保证接收区域，再比较其移动时间和最坏后验定位半径。

\subsection{问题三分析}
全向源的无信号观测能够排除检测点周围半径 $1000\,\mathrm{m}$ 的区域。因而可以设计覆盖整个目标圆的检测点集合，既用于发现存在的频道，也用于证明其余频道不存在；在此基础上把搜索、顺路补测和安全清除交织执行。

\subsection{问题四分析}
定向源使无信号同时可能表示距离过远或位于发射方向背面，问题三的圆盘排除不再成立。需要构造同时满足距离覆盖和方向覆盖的检测点布局，并仅用正观测更新位置角域。
